\section{Sentiment Analysis}
\textit{Design and implement a PyTorch pipeline to train a Transformer model for a text
classification task (e.g., sentiment analysis). Include dataset preprocessing, model
definition, training loop, and evaluation. Discuss challenges such as overfitting and
computational cost.}

Sentiment analysis is a supervised text classification task
in which the goal is to assign a discrete sentiment label,
such as positive or negative, to a given input sentence.
Transformer-based models are particularly well-suited for
this task, as self-attention enables the model to capture
contextual dependencies across an entire sequence.
In this section, we design and implement a PyTorch pipeline
for training a Transformer encoder on a sentiment
classification dataset, following the implementation
outlined in~\cite{HowUsePyTorch}.

The pipeline begins with dataset preprocessing.
Raw text sentences are first tokenized using a simple
whitespace-based tokenizer, and a vocabulary is constructed
from the training corpus.
Each token is mapped to a unique index, allowing sentences
to be converted into sequences of integers.
Since Transformer models operate on fixed-size batches,
all sequences are padded to a common length.
A corresponding padding mask is created to ensure that
padded tokens do not contribute to the attention scores
or pooled representations during training.

The model architecture consists of an embedding layer
followed by positional encoding and a Transformer encoder.
The embedding layer maps token indices to dense vectors,
while positional encoding injects word-order information
that is otherwise absent from the attention mechanism.
These embeddings are then passed through multiple layers
of a Transformer encoder with multi-head self-attention,
allowing the model to learn contextualized representations
of each token in the input sequence.

To obtain a fixed-size representation suitable for
classification, token-level embeddings are aggregated
using mask-aware mean pooling.
Specifically, embeddings corresponding to padded positions
are excluded from the average, ensuring that sequence length
variation does not bias the representation.
The pooled sentence representation is then passed through
a linear classifier, which outputs logits corresponding
to the sentiment classes.

The model is trained using the cross-entropy loss function,
which is appropriate for multi-class classification tasks.
Optimization is performed using the AdamW optimizer, which
combines adaptive learning rates with decoupled weight
decay to improve generalization.
The training loop iterates over mini-batches of data,
computes forward passes through the model, backpropagates
the loss, and updates the model parameters.
Performance is evaluated periodically on a validation set
to monitor convergence and detect overfitting.

Several practical challenges arise when training Transformer
models for sentiment analysis.
Overfitting is a common issue, particularly when training
on small or moderately sized datasets.
This is mitigated through the use of dropout within the
Transformer layers and weight decay in the optimizer.
Another important consideration is computational cost.
Self-attention has quadratic complexity with respect to
sequence length, which can lead to increased memory usage
and slower training for long inputs.
As a result, careful selection of maximum sequence length
and batch size is necessary to balance performance and
efficiency.

The full implementation of this sentiment analysis pipeline,
including the data preprocessing
is provided in \mintinline{text}{code/sentiment_analysis.py}.
